\subsection{基本対称式型}
導入に続いて,先程の漸化式を考えていこう.
\[
  a_1=3, \quad a_{n+1}=a_n^2-2
\]
この漸化式を解くに当たって邪魔をしているのは$-2$の項である.
これをうまく消せるような操作を考えたい.
また,\,2乗があることからうまく因数分解できないか考えてみよう.

以上の$2$と打ち消し会えるように,\,$2$が生まれるような対称式を考える.

2乗の項は正だから,
\[
  (a+b)^2=a^2+2ab+b^2
\]
で言うところの$2ab$が$2$になってくれるように調整すればうまくいきそうだ.
そのようなふうに考えていくと,
\[
  \left(\alpha+\frac{1}{\alpha}\right)^2=\alpha^2+2+\frac{1}{\alpha^2}
\]
という式が思い当たる.
実際,ここに$-2$をすると
\[
  \alpha^2+\frac{1}{\alpha^2}
\]
が残るから,成り立ちそうなことが分かる.
この実験を元に解答に移る.

\noindent \textbf{【問題】}
\begin{quote}
  次の漸化式で定義される数列の一般項を求めよ.
  \[
    a_1=3, \quad a_{n+1}=a_n^2-2
  \]
\end{quote}
\textbf{【解答】}\\一般項が
  \[
    b_n=\alpha^{2^{n-1}}+\alpha^{-2^{n-1}}
  \]
  で表される数列$b_n$を考える.
  \begin{align*}
    b_{n+1}&=\alpha^{2^{n}}-\alpha^{-2^{n}}\\
    &=\left(\alpha^{2^{n-1}}\right)^2-\left(\alpha^{-2^{n-1}}\right)^2\\
    &=\left(\alpha^{2^{n-1}}+\alpha^{-2^{n-1}}\right)^2-2\\
    &=b_n^2-2
  \end{align*}
  これは与式の漸化式と一致する.また,
  \[
    b_1=a_1=\alpha+\alpha^{-1}
  \]
  となる$\alpha$を考えると,
  \begin{align*}
    &\alpha+\alpha^{-1}=3\\
    &\alpha^2-3\alpha+1=0\\
    &\therefore \quad \alpha=\frac{3\pm \sqrt{5}}{2}
  \end{align*}
  得られた$\alpha$の2解は逆数の関係にある.
  以上より,
  \[
    a_n=\left(\frac{3+ \sqrt{5}}{2}\right)^{2^{n-1}}+\left(\frac{3- \sqrt{5}}{2}\right)^{2^{n-1}}
  \]
\begin{quote}
  
\end{quote}

ここで,定数を $-2$ から $-1$ に変えてみよう.
\[
  a_{n+1}=a_n^2-1
\]
同じ置換を試すと,
\begin{align*}
  a_n^2-1
    &=\left(u+u^{-1}\right)^2-1\\
    &=u^2+u^{-2}+1,
\end{align*}
となる.しかし,次の項として欲しい形は $u^2+u^{-2}$ である.
余分な $1$ が残るため,同じ対称式の族の中に戻ってこない.

漸化式の $-2$ と $-1$ は,見た目には小さな違いである.
しかし,対称式を保つ恒等式にとっては,その違いが決定的になる.
「ほとんど同じ式だから,ほとんど同じ方法で解ける」とは限らない.
解法が成立するかどうかは,係数の大きさではなく,恒等式が完全に一致
するかどうかで決まる.

同じ考え方は,三次式にも現れる.
\[
  \left(\alpha+\frac{1}{\alpha}\right)^3
  -3\left(\alpha+\frac{1}{\alpha}\right)
  =\alpha^3+\frac{1}{\alpha^3}
\]
このように,対称な組み合わせを選ぶと,べき乗の複雑な展開が
再び同じ種類の式へ戻る.ただし,恒等式が少しでもずれると,その閉じた
構造は失われる.これが基本対称式型の面白さである.

これを発想にもつ漸化式を例題としよう.

\noindent\textbf{【問題】}\\
$\displaystyle a_1=\frac{8}{3},\quad a_{n+1}=a_n^3+3a_n$

\noindent\textbf{【解答】}\\
$a_n=b_n-\frac{1}{b_n}$\\
と置くと.

$\displaystyle \alpha-\alpha^{-1}=\frac{8}{3}$
を解くと,
$\displaystyle \alpha=-\frac{1}{3},\,3$
